\documentclass[11pt]{article}

\usepackage[margin=1in]{geometry}
\usepackage{amsmath,amssymb,amsthm,mathtools}
\usepackage{enumitem}
\usepackage{hyperref}
\hypersetup{colorlinks=true,linkcolor=blue,citecolor=blue,urlcolor=blue}
\usepackage{orcidlink}
\usepackage{fontawesome5}

\newcommand{\ZZ}{\mathbb{Z}}
\newcommand{\NN}{\mathbb{N}}

\title{Scope Honesty in the Collatz-Adjacent Literature:\\
A Running Calibration Against Complete-Proof Claims}

\author{Renato Augusto Tavares~{\small\href{mailto:dr.renatotavares@gmail.com}{\textcolor{black}{\faEnvelope}}}~\orcidlink{0009-0002-0196-3311}\\
\normalsize Universidade Federal de Goiás}
\date{\today}

\begin{document}

\maketitle

\begin{abstract}
A companion note to this survey catalogs twelve claims of a complete
proof or disproof of the Collatz conjecture, all of them invalid. This
note turns to a different population: sixteen papers, encountered in
the same broader literature survey, that make no such claim --- results
about Collatz-adjacent structures, heuristics, or partial statistics.
Read in full and checked computationally where possible, this
population behaves very differently from the proof-claim set. The
distinguishing axis is not error rate but scope honesty: with no known
exception in this sample, these papers correctly label a conjecture a
conjecture and a theorem a theorem, in sharp contrast to the systematic
overclaiming that defines the companion catalog. Genuine errors do
occur --- most notably two in a seven-paper program by a single
prolific author, one
in an otherwise careful standalone result --- but they cluster
differently: the prolific author's errors are both the same shape (a
conjectural or preliminary claim inconsistent with an
already-established statement elsewhere in the same text, not an
arithmetic slip) and contained to a single section each, while the
standalone result's error sits in its titular theorem. This note is
explicitly a partial, running document: it consolidates sixteen items
reviewed in depth so far, out of a much larger surveyed collection, and
is structured to grow as more items are processed.
\end{abstract}

\section{Introduction}
\label{sec:intro}

A companion note, \emph{Errors in Claimed Complete Proofs (and
Disproofs) of the Collatz Conjecture}, catalogs twelve claims of a
complete solution to the conjecture, encountered during a broader
literature survey, and shows that all twelve fail --- most by the same
handful of recurring errors. That population is unusual: every item in
it explicitly claims to settle the conjecture, a strong claim that
invites strong scrutiny.

This note is about the rest of the survey: papers that study
Collatz-adjacent structures --- accelerated-dynamics reformulations,
statistical models of orbit behavior, combinatorial properties of
parity sequences, coordinate systems for the state space --- without
claiming to prove or disprove the conjecture itself. This is a much
larger and more heterogeneous population than the proof-claim set, and
this note does not attempt to cover all of it: as of this version, it
consolidates sixteen items that have been read in full and reviewed in
depth, out of a considerably larger set encountered in the survey. It
is written as a running document, expected to grow as more of the
surveyed collection is processed in later sessions, not as a finished
census.

The question this note asks of the sixteen items it covers is not
primarily ``are they correct'' --- most are, and correctness in
elementary or standard mathematics applied competently is not
surprising. The more informative question, and the one that lets this
note function as a calibration point against the companion catalog, is
scope honesty: does the paper claim only what it establishes? Section
\ref{sec:rc} examines the largest coherent block in this sample, a
seven-paper program by a single author applying one classical
ergodic-theory concept per paper to the same underlying quantity.
Section \ref{sec:others} covers the remaining nine items, independent
of one another and of the program in Section \ref{sec:rc}. Section
\ref{sec:discussion} draws out the contrast with the companion catalog
and states what containment of error does and does not generalize to.
Section \ref{sec:conclusion} concludes.

\section{Methodology}
\label{sec:methodology}

Each item was read in full and its concrete, checkable claims ---
identities, closed-form formulas, theorems with an explicit statement
--- were independently verified computationally wherever they admit a
numerical test. Where a paper distinguishes proved results from
conjectural ones, that labeling was checked against what the paper's
own proof actually establishes, not taken at face value. In the
handful of cases where the judgment involved was subtle, the analysis
was cross-checked against an independent large-language-model reviewer.
As in the companion catalog, every specific claim made below about a
named paper traces to this review process, not to a paraphrase of it;
full technical detail for each item was recorded during review in
internal working notes and is summarized, not reproduced, here.

\section{The Ruiz Castillo program}
\label{sec:rc}

Seven papers by Juan Carlos Ruiz Castillo (ScD., Universidad de San
Carlos de Guatemala) appear in this sample, each applying one classical
concept from ergodic theory or large-deviations theory --- pressure,
a residual ``geometry,'' dissipative bounds, a central limit theorem,
a transfer operator, Gibbs measures, a variational principle, a
large-deviations rate function --- to the same elementary quantity:
the residual drift $L_k(n) = k\log_2 3 - A_k(n)$, where $A_k(n)$ is the
sum of the first $k$ 2-adic valuations along the accelerated Collatz
map $U(n) = (3n+1)/2^{v_2(3n+1)}$ --- the standard cycle-equation drift
already used elsewhere in this research program. None of the seven
claims to prove the Collatz conjecture; each says so explicitly. The
presentation is heavily repetitive across the series (recurring
decorative call-out boxes restating the same conceptual chain,
self-citation rates of 75--100\% revealing on the order of 19--20
near-identical titles by the same author), but the underlying formal
rigor is notably stable: five of the seven papers contain no error we
could find, and the two genuine errors found are both contained to a
single section and do not affect the papers' own later conclusions.

\textbf{Geometría Residual} \cite{RuizCastillo2026a} applies standard
thermodynamic formalism (pressure as a supremum over invariant
measures, a Fisher-metric second derivative) to the residual drift,
correctly and without claiming to prove Collatz. The same trivial fact
--- convexity implies a nonnegative second derivative --- is presented
as five separately named results across the paper, framed by ten
decorative call-out boxes restating the same conceptual chain, with no
numerical content anywhere (both figures are explicitly labeled
``conceptual'') and 100\% self-citation (13 of 13 references). We
computed explicitly, in closed form, the quantity the paper leaves
abstract: under the standard i.i.d.\ model already used in this
research program, $\Lambda(t) = t(1-\log_2 3) - \log(2-e^t)$, whose
second derivative diverges as $t \to \log 2^-$ --- a genuine, if
incidental, feature the paper never computes because it never derives
$\Lambda(t)$ explicitly.

\textbf{Dissipative Bounds} \cite{RuizCastillo2026b}, the same author's
second paper in this sample, repeats the pattern: correct elementary
identities (an exact affine unrolling of the accelerated map, a
universal dissipation bound), no numerical content, and self-citation
at 75\% (12 of 16 references), with the self-citation list itself
revealing roughly a dozen further near-identical titles by the same
author on the same ``residual decomposition.''

\textbf{Teorema Central del Límite Residual} \cite{RuizCastillo2026c}
labels its central claim honestly: the title says ``Teorema,'' but the
result is stated in the body as ``Conjetura 4.2,'' with five explicit
unestablished hypotheses (existence of a residual Gibbs measure,
ergodicity, a suitable function space, a spectral gap, positive
residual variance), and the paper's own conclusion states plainly that
this framework does not prove the Collatz conjecture. We tested the
conjecture's observable prediction --- asymptotic normality of the
residual drift's fluctuations --- directly on real Collatz orbits
(not the abstract i.i.d.\ model): variance stabilizes near 2, skewness
toward 0, and kurtosis toward 3 as the window length grows, a result
the paper itself never tests numerically.

\textbf{Operador de Transferencia Residual} \cite{RuizCastillo2026d}
contains this program's first genuine error. Its Proposición 5.3
claims $\lim_{t\to\infty} L_t(1) = 0$; its own proof derives the closed
form $L_t(1) = e^{(\log_2 3 - 1)t}/(1-e^{-t})$ and correctly observes
that this ``grows exponentially as $t\to\infty$,'' immediately followed
by the proof's closing mark with the ``$=0$'' claim never corrected. We
confirmed numerically that $L_t(1)$ grows from 2.84 at $t=1$ to
roughly $6.4\times 10^{50}$ at $t=200$, strictly increasing throughout
--- the opposite of the stated limit. This is a statement-versus-proof
inconsistency, not a computational error: the proof itself derives the
correct asymptotic and the surrounding text uses that correct
asymptotic to motivate the next section's normalization. The error is
confined to a preliminary calculation in Section 5; no later,
explicitly conjectural result (Sections 6--8, on the restricted
transfer operator) depends on it.

\textbf{Medidas de Gibbs Residuales} \cite{RuizCastillo2026e} and
\textbf{Principio Variacional Residual} \cite{RuizCastillo2026f} return
to the error-free pattern: correct restatements of the same core
identity, classical convex-analysis apparatus (suprema of affine
functions, Legendre--Fenchel duality) correctly applied, and every open
result honestly labeled ``Conjetura,'' never ``Teorema'' or
``Proposición.''

\textbf{Grandes Desviaciones Residuales} \cite{RuizCastillo2026g}
contains this program's second genuine error, structurally similar in
kind to the first but located differently: its Definición 3.1 defines
a rate function $I_{\mathrm{RC}}(x)$ via the one-sided tail event
$\{L_k/k \ge x\}$, and its already-proved Proposición 3.4 (Section 3)
correctly shows this one-sided rate function is nondecreasing. Its
conjectural Section 7, however --- Figura 1, Conjetura 7.3, and
Conjetura 7.5, mutually consistent with one another --- describes a
two-sided, classical Cramér rate function, positive on both sides of
its unique zero, contradicting the monotonicity already established in
Section 3. We confirmed the contradiction by three independent routes
(the standard one-sided large-deviations formula, Monte Carlo
simulation, and an exact small-$k$ negative-binomial calculation), all
agreeing that the correct one-sided rate is identically zero below the
typical drift, not merely at one point. As with the previous case, the
error is a conjectural section's inconsistency with an already-proved
proposition earlier in the same text, not an arithmetic mistake, and is
confined to Section 7; the paper's conclusion (Section 8) does not rely
on the specific form of $I_{\mathrm{RC}}$.

\section{Nine independent papers}
\label{sec:others}

The remaining nine items in this sample share no authorship and little
subject matter beyond Collatz-adjacent structures.

\textbf{Adnan \& Dar} \cite{AdnanDar2026} extend the $3n+1$ rule to
decimals in $(0,1)$ by defining ``parity'' via a number's last
significant decimal digit, concluding universal divergence. All worked
examples in the paper check out arithmetically, but the construction is
conceptually ill-posed: ``last significant digit'' is undefined for any
decimal with an infinite expansion (a measure-zero exception in the
naive reading, but one that excludes $1/3$, $1/7$, and every
irrational), and even restricted to finite decimals the divergence is a
trivial artifact of an unnormalized $+1$ term, not a structural analogue
of the real Collatz mechanism (where $3n+1$ is always even for odd $n$
by construction). This is not a proof claim about the original
conjecture and does not belong in the companion catalog.

\textbf{Gilbert} \cite{Gilbert2026} constructs an explicit conjugation
between the pruned Collatz graph and the nonzero integers, encoding the
discarded residue class mod 3 via sign, explicitly disclaiming any
proof of the conjecture. Every theorem we checked --- the conjugation
itself, pruning of multiples of 3, the accelerated map, the parent-count
formula, and a connection to OEIS sequence A254046 --- holds exactly.
No error found.

\textbf{Seymour (Steiner Sentence Length)} \cite{Seymour2026a} derives a
matrix of one-step transition probabilities among residues mod 8
(Theorem 2.1), which we confirmed correct, and uses it to derive a
claimed closed-form distribution for ``Steiner sentence'' length,
$3^{k-1}/4^k$ (Theorem 5.1), rejecting the naive alternative $(1/2)^k$
via a reported $\chi^2$ test. Our independent reimplementation of the
paper's own definitions gives the opposite result: simulation matches
$(1/2)^k$, and an exact arithmetic counterexample ($n=68567$) shows a
sentence of length 1 whose entry residue is 7, a case the paper's proof
of Theorem 5.1 never accounts for (it only counts sentences whose
\emph{first} word already closes, silently missing words with entry 3
or 7 that can close on their first step). We report this with real
caution: the paper claims a complete Lean 4/Mathlib formalization with
no \texttt{sorry}, and we do not have access to that formalization
file; it is possible the Lean proof correctly verifies a definition
that diverges subtly from the prose statement, in which case the gap
would be in the prose-to-formalization translation rather than in the
formal proof itself. Based strictly on the definitions as written in
the paper, however, both an exact counterexample and a large-scale
simulation disagree with the central result.

\textbf{Seymour (Regular Expression Language)} \cite{Seymour2026b}, an
explicitly labeled working paper by the same author, separates proved
theorems from conjectures throughout. Two central results (Proposición
3.1 on Steiner-circuit arithmetic, Teorema 3.6 on run lengths) checked
out exactly in large-scale testing. One small, contained error: Corolario
2.2 claims an exit residue mod 24 depends on the parity of an auxiliary
variable $t$; testing shows the claimed value holds only when $t$ is
even, and the true residue is the same constant regardless of $t$'s
parity, contradicting the paper's own claim that the map $t \mapsto
(3t+1)/2$ preserves parity (false in general). This corollary is not
listed among the paper's own summary of proved results, and the error
does not affect the working paper's central conjecture, which follows
almost immediately from the already-verified Theorem 2.1 of the
companion paper.

\textbf{Anthony} \cite{Anthony2026} builds an elaborate formal apparatus
(Riemann P-equation, monodromy, hypergeometric identities) around a
two-field reformulation of Collatz dynamics via the Möbius coordinate
$\Phi(n) = 1/(n+1)$. Every checkable claim --- the Möbius
reformulation, a theorem excluding indefinite persistence of the
$v_2(3n+1)=1$ regime, and several standard digamma-function identities
--- checked out exactly. The paper is unusually careful throughout to
separate proved theorems from structural analogies (explicitly: ``this
is an analogy of form, not of kind'') and from conditional results
whose conditions are not established, and it corrects, within its own
text, an earlier erroneous claim from a prior version of the same
research line.

\textbf{Hikawa} \cite{Hikawa2026} studies finite combinatorial
structures of parity vectors, independent of whether the Collatz
conjecture is true or false --- a caveat the paper repeats in nearly
every section. Both central theorems we tested (an explicit bijection
between vector classes, and a rigidity property of a correction term's
2-adic valuation) held exactly by brute force, including the paper's
own stated exception case.

\textbf{Olgac} \cite{Olgac2026} is a three-page note connecting two
earlier, uncataloged claims by the same author, one about Collatz and
one claiming a complete proof of the Goldbach conjecture (outside this
survey's scope). Its only Collatz-specific content is a tautology
(root-disjoint branches of the reverse tree, an immediate consequence
of unique factorization) rather than a genuine result; the remainder of
the note is not stated with enough mathematical precision to evaluate
as correct or incorrect.

\textbf{Williams} \cite{Williams2026} organizes the positive integers by
their 3-smooth factorization and studies Collatz dynamics as a
deterministic diagonal flow within that organization. Every theorem we
tested --- a diagonal-dynamics identity, a theorem on prime-free rows,
a bijective partition, an asymptotic count of admissible positions ---
held exactly. Two small, self-contained errors in OEIS citations were
found (one sequence number that does not match the cited property, one
off-by-one description), neither affecting any proof in the paper. The
paper separates proved theorems from computational observations and
explicit open problems throughout, and discloses its use of generative
AI (for drafting and formatting, with declared full human
verification).

\textbf{Fu, Liu \& Wang} \cite{FuLiuWang2026} propose a Poisson-process
mechanism explaining a previously observed Gamma-type statistic for
Collatz orbit lengths. Every checkable claim --- the expectation and
variance of the 2-adic valuation, a closed-form approximation linking
orbit length to a starting value, and the algebra of a cycle-closure
condition --- held exactly. The paper is explicit, correctly and in
contrast to the flawed proof attempt reviewed in the companion catalog
under the same general argument style, that its asymptotic
cycle-closure obstruction is not a proof that no nontrivial cycle
exists.

\section{Discussion}
\label{sec:discussion}

The contrast that holds across this entire sample, with no known
exception, is scope honesty: not one of these sixteen papers dresses a
conjecture as a proof of the Collatz conjecture itself, even where the
underlying mathematics is elementary and heavily repeated (the Ruiz
Castillo program) and even where a central result is probably wrong
(Seymour's Steiner sentence paper). This is the axis on which this
population differs from the companion catalog's twelve cases, where
every single item makes the strong claim and every single item fails to
support it. Error rate is not the distinguishing feature here --- both
populations contain real errors --- overclaiming is.

Containment of error is a narrower, true observation, but it holds
about the Ruiz Castillo program specifically, not as a property of this
population as a whole. Both of that program's genuine errors are the
same shape (a conjectural or preliminary claim inconsistent with an
already-established statement elsewhere in the same text, not an
arithmetic slip) and both are confined to a single section, leaving the
paper's later, explicitly conjectural material untouched. Seymour's
Steiner sentence paper is the sharp counterpoint: its error sits in the
titular Theorem 5.1, the paper's central claimed result, not a
peripheral lemma. The interesting fact is not that errors in this
literature are contained --- they are not, in general --- but that the
most prolific single-author program in this sample errs in a
consistently contained way, while the sharpest standalone result in the
sample errs at its center. That difference is a fact about how
particular authors err, not a generalization about the population.

We note, finally, that this chapter's coverage is partial by
construction: sixteen items reviewed in depth here, drawn from a
considerably larger surveyed collection that includes both further
unprocessed Collatz-adjacent papers and a separate body of established
reference literature (Tao's almost-all-orbits theorem, the
Kontorovich--Lagarias stochastic model, and similar) that is cited
throughout this research program but was never a target of this kind of
review. Future revisions of this note are expected to add further
items as they are processed; the structure above --- one section per
coherent author or program, individual paragraphs per independent item
--- is meant to accommodate that growth without requiring
reorganization.

\section{Conclusion}
\label{sec:conclusion}

Sixteen Collatz-adjacent papers, none claiming to prove or disprove the
conjecture, reviewed in full: nine contain no error we could find; four
contain a small, contained error (two in the Ruiz Castillo program, one
in a Seymour working paper, and one, a pair of citation-only slips, in
Williams); one --- Seymour's Steiner sentence paper --- has a likely
error in its central, titular theorem rather than a peripheral one; and
two make no claim evaluable as straightforwardly correct or incorrect
(Adnan \& Dar's construction is conceptually ill-defined; Olgac's only
Collatz-specific content is tautological). Across all sixteen, the
theorem-versus-conjecture distinction is maintained with a rigor this
survey's
companion catalog of proof claims never shows. That contrast --- scope
honesty, not error rate --- is this note's calibration point, and it is
expected to sharpen, not change in kind, as more of the surveyed
collection is added in later revisions.

\section*{Acknowledgments}

Parts of the verification described in this note (case reconstruction,
computational checks, and cross-checking of subtle mathematical
judgment calls) were carried out with the assistance of Claude (models
in the Claude 4 and 5 families, Anthropic), used as a research
assistant under the direction and review of the author.

\begin{thebibliography}{99}

\bibitem{RuizCastillo2026a} J.\ C.\ Ruiz Castillo, \emph{Geometría
  Residual de Ruiz Castillo para la dinámica acelerada de la Conjetura
  de Collatz}, Zenodo, DOI 10.5281/zenodo.21271452, 2026.

\bibitem{RuizCastillo2026b} J.\ C.\ Ruiz Castillo, \emph{Dissipative
  Bounds and Ruiz Castillo Residual Decomposition in the Accelerated
  Dynamics of the Collatz Conjecture}, ResearchGate preprint.

\bibitem{RuizCastillo2026c} J.\ C.\ Ruiz Castillo, \emph{Teorema Central
  del Límite Residual de Ruiz Castillo para la dinámica acelerada de la
  Conjetura de Collatz}, ResearchGate preprint.

\bibitem{RuizCastillo2026d} J.\ C.\ Ruiz Castillo, \emph{Operador de
  Transferencia Residual de Ruiz Castillo y estructura espectral en la
  dinámica acelerada de la Conjetura de Collatz}, ResearchGate preprint.

\bibitem{RuizCastillo2026e} J.\ C.\ Ruiz Castillo, \emph{Medidas de
  Gibbs Residuales de Ruiz Castillo y estados de equilibrio en la
  dinámica acelerada de la Conjetura de Collatz}, ResearchGate preprint.

\bibitem{RuizCastillo2026f} J.\ C.\ Ruiz Castillo, \emph{Principio
  Variacional Residual de Ruiz Castillo y formalismo termodinámico para
  la dinámica acelerada de la Conjetura de Collatz}, ResearchGate
  preprint, 2026.

\bibitem{RuizCastillo2026g} J.\ C.\ Ruiz Castillo, \emph{Grandes
  Desviaciones Residuales de Ruiz Castillo en la dinámica acelerada de
  la Conjetura de Collatz}, Zenodo, DOI 10.5281/zenodo.20767811, 2026.

\bibitem{AdnanDar2026} J.\ Adnan and S.\ A.\ Dar, \emph{Behavior of a
  Decimal-Parity-Based $3n+1$ Mapping}, Global Journal of Pure and
  Applied Mathematics, 2026.

\bibitem{Gilbert2026} J.\ S.\ Gilbert, \emph{A Collatz-Equivalent Map on
  the Nonzero Integers}, Preprints.org, DOI
  10.20944/preprints202607.0575.v1, 2026.

\bibitem{Seymour2026a} J.\ Seymour, \emph{First-Principles Derivation of
  the Steiner Sentence Length Distribution}, wildducktheories.com
  preprint, 2026.

\bibitem{Seymour2026b} J.\ Seymour, \emph{A Regular Expression Language
  for the Collatz Graph}, wildducktheories.com working paper, 2026.

\bibitem{Anthony2026} M.\ M.\ Anthony, \emph{A Two-Field Propagation
  Model for the Collatz Map: A Möbius Reformulation, Monodromy
  Structure, and a 2-adic Non-Expansion Theorem}, ResearchGate
  preprint, 2026.

\bibitem{Hikawa2026} K.\ Hikawa, \emph{Finite-Dimensional Combinatorial
  and Arithmetic Structures of Parity Vectors for the Accelerated
  Collatz Map}, ResearchGate preprint.

\bibitem{Olgac2026} E.\ Olgac, \emph{Technical Note: Structural Dualism
  in Integer Architectures}, ResearchGate preprint.

\bibitem{Williams2026} J.\ Williams, \emph{A Coordinate System for
  Collatz Dynamics}, arXiv:2607.01718, 2026.

\bibitem{FuLiuWang2026} W.\ Fu, X.\ Liu, and Y.\ Wang, \emph{Emergence of
  Gamma-Type Upward-Phase Statistics in the Collatz Map: An Effective
  Poisson Process Mechanism}, arXiv:2606.26811, 2026.

\end{thebibliography}

\end{document}
